\section{The structure-theory programme}\label{sec:programme}

The preceding sections produced, from an almost-idempotent positive map, an
$\eps$-JB algebra: a finite-dimensional real space with a commutative product
$a\bullet b := P(a\jp b)$, a unit $\Id$, and an order-unit norm, satisfying the
JB axioms up to an error of order $\eps$ (see \Cref{sec:epsjb}). The goal of the
structure-theory programme is to remove the error entirely: to show that such an
object is, up to a controlled deformation, a genuine JB-algebra. This section
states that goal honestly as an open target, records the proof strategy, and
isolates the one quantitative obstruction together with candidate routes around
it.
Nothing below is claimed as a completed proof; every named target carries the
tag \status{(open)}.

\subsection{The target}

The model to imitate is Kitaev's structure theorem for approximate
$C^{*}$-algebras. In his terminology an \emph{\eps-$C^{*}$ algebra} is a
finite-dimensional space with a product, involution and norm satisfying the
$C^{*}$ axioms up to \eps, and a \emph{$\delta$-isomorphism} is a linear
bijection preserving unit and product up to a relative error $\delta$. His main
theorem reads, in source notation after harmless macro expansion:
\begin{quote}
``Any finite-dimensional $\eps$-$C^{*}$ algebra $\mathcal{A}$ is
$O(\eps)$-isomorphic to some $C^{*}$ algebra $\mathcal{B}$. (Importantly, the
implicit constant in $O(\eps)$ does not depend on $\mathcal{A}$ or its
dimensionality.)''
\end{quote}
\cite{Kitaev2024}, Theorem \texttt{th\_main}. The decisive feature is the last
sentence: the constant is \emph{dimension-free}. The Jordan analogue we are
after is the following.

% PROV: Kitaev2024 approximate_algebras.tex:460-462 (th_main) -- C*-statement; ours is the declared Jordan analogue (open).
\begin{openproblem}[Jordan structure theorem]\label{thm:jordan-structure}
\status{(open)}
There are universal constants $C,\eps_{0}>0$ \emph{(independent of dimension)}
such that for every $\eps<\eps_{0}$ and every finite-dimensional $\eps$-JB algebra
$\mathcal{A}$ there exist a genuine finite-dimensional JB-algebra $\mathcal{B}$
and a linear bijection $v\colon\mathcal{B}\to\mathcal{A}$ that is a
$C\eps$-Jordan-isomorphism, i.e.
\[
   \norm{v(\Id_{\mathcal{B}})-\Id_{\mathcal{A}}}\le C\eps,\qquad
   \norm{v(b\bullet b')-v(b)\bullet v(b')}\le C\eps\,\norm{b}\,\norm{b'},
\]
and $(1-C\eps)\norm{b}\le\norm{v(b)}\le(1+C\eps)\norm{b}$ for all
$b,b'\in\mathcal{B}$.
\end{openproblem}

This is the analogue of Kitaev's statement with the associative product replaced
by the Jordan product $\bullet$ and ``$C^{*}$ algebra'' replaced by
``JB-algebra''. As of the present report it is a target, not a theorem. What we
\emph{can} state is the strategy and the precise point at which the
dimension-free constant is in question.

\subsection{The strategy}

Kitaev's proof builds the model algebra $\mathcal{B}$ incrementally and controls
the error at each step by a cohomological argument. The Jordan programme mirrors
the four moving parts.

\paragraph{Vocabulary for the programme.}
An \emph{idempotent element} is an element $e$ with $e\bullet e=e$. A
\emph{Jordan frame} is a finite list of nonzero idempotents that are pairwise
orthogonal, $e_i\bullet e_j=0$ for $i\ne j$, and sum to the unit; it plays the
role of diagonal matrix units. The \emph{Peirce decomposition} is the splitting
of the algebra into the eigenspaces of multiplication by an idempotent. A
\emph{cochain} is a multilinear correction term; a \emph{cocycle} is a defect
that satisfies the compatibility identity forced by the Jordan identity; a
\emph{coboundary} is a defect removable by changing the comparison map. A
\emph{contracting homotopy} or \emph{splitting} is a bounded linear recipe that
turns a cocycle into such a correction. The \emph{projective tensor norm} records
the cost of writing a tensor as a sum of simple tensors; dimension-free control
means this cost is bounded by a universal constant.

\medskip
\noindent\textbf{(a) Approximate idempotents and a Jordan frame.}
A genuine JB-algebra is coordinatized by a complete system of orthogonal
idempotents (a \emph{Jordan frame}). The first step produces such idempotents
\emph{approximately}: elements $e=e^{*}$ with $\norm{e\bullet e-e}$ small and
not too close to $0$ or $\Id$. In the associative case these are the nontrivial
projections produced by a fixed-point argument \cite{Kitaev2024}; the spectral
idempotent $P$ of \Cref{sec:spectral} is the map-level device that produced the
$\eps$-JB product. It does \emph{not} by itself produce idempotent elements in
$A$; constructing approximate Jordan-frame elements is a separate open step. The
Jordan frame would then be refined by repeatedly splitting an approximate
idempotent into two, exactly as the associative model algebra is grown from $\C$
to $\C\oplus\C$ and onward.

\smallskip\noindent\textbf{Agent B review correction.}
The earlier draft blurred the idempotent linear map $P$ with idempotent elements
of the approximate algebra. Layer 1 still needs its own approximate-frame
construction.

\medskip
\noindent\textbf{(b) Approximate Peirce decomposition.}
A genuine idempotent $e$ in a Jordan algebra induces the Peirce decomposition
$\mathcal{A}=\mathcal{A}_{1}\oplus\mathcal{A}_{1/2}\oplus\mathcal{A}_{0}$ into
the $1$, $\tfrac12$, $0$ eigenspaces of the multiplication operator
$L_{e}\colon a\mapsto e\bullet a$. For an approximate idempotent the operator
$L_{e}$ has spectrum clustered near $\{0,\tfrac12,1\}$, and the associated
spectral projections give an \emph{approximate} Peirce decomposition with
$O(\eps)$ cross-terms. This decomposes $\mathcal{A}$ into pieces on which the
product is closer to the model multiplication rules.

\medskip
\noindent\textbf{(c) Approximate coordinatization.}
On the joint Peirce spaces of a frame of rank $\ge 3$, the Jordan
coordinatization theorem (the Jordan--von Neumann--Wigner coordinatization of
$\mathrm{H}_{n}(\mathfrak{C})$, \cite{JordanVonNeumannWigner1934}) identifies
the algebra with hermitian matrices over a coordinate $*$-algebra. The
approximate version reads off coordinate elements from the approximate Peirce
pieces and checks the coordinate relations up to $O(\eps)$. This step delivers a
candidate genuine JB-algebra $\mathcal{B}$ together with a linear inclusion
$v\colon\mathcal{B}\to\mathcal{A}$ that is multiplicative up to a (possibly
large) error $\delta$.

\medskip
\noindent\textbf{(d) Error reduction.}
The inclusion $v$ from (c) is only a $\delta$-inclusion with $\delta\gg\eps$.
The final step replaces $v$ by a nearby map whose defect is $O(\eps)$,
\emph{uniformly in the dimension}. This is the cohomological heart of the
argument and is treated separately below.

The algebraic inputs to (d) are classical and we cite them as facts. For a
finite-dimensional Jordan algebra $J$ of characteristic $0$ and a $J$-module
$M$, the two Whitehead lemmas assert the vanishing of the first two cohomology
groups in the semisimple/separable case.

% PROV: ChuRusso2016 (Thm 1.1, 1st Whitehead lemma, attrib. Jacobson) via agent-A/notes/ingestion-results-2026-06-01.json -- H^1=0; Penico1951 (2nd Whitehead lemma) via same -- H^2=0. Extraction-level provenance; primary PDF text still pending independent byte-check.
\begin{theorem}[Whitehead lemmas for Jordan algebras, \status{(source-extracted; pending primary byte-check)}]\label{thm:whitehead}
Let $J$ be a finite-dimensional semisimple Jordan algebra over a field of
characteristic $0$ and let $M$ be a finite-dimensional $J$-module.
\begin{enumerate}
\item \emph{(First Whitehead lemma.)} Every derivation $f\colon J\to M$ is
inner; equivalently $H^{1}(J,M)=0$.
\item \emph{(Second Whitehead lemma.)} If moreover $J$ is separable, every
Jordan $2$-cocycle is a coboundary; equivalently $H^{2}(J,M)=0$.
\end{enumerate}
The first lemma is due to Jacobson; the second to Penico, who used it to prove
the Wedderburn principal theorem for Jordan algebras.
\end{theorem}

\noindent The first part is the Jordan analogue of the first Whitehead lemma,
recorded as Theorem~1.1 of \cite{ChuRusso2016}: ``Every linear $f\colon J\to M$
with $f(ab)=f(a)b+af(b)$ is inner.'' The second part is the Jordan analogue of
the second Whitehead lemma, Theorem~1.2 of \cite{ChuRusso2016}, attributed
there to \cite{Penico1951}: ``Every Jordan $2$-cocycle $f$ is a coboundary.''
The vanishing of $H^{1}$ and $H^{2}$ is what would make the exact correction
step possible \emph{in principle}; the issue is the size of the correction.

The vanishing has a concrete mechanism, which is the second cited input. The
\emph{multiplication algebra} $M(J)\subseteq\operatorname{End}(J)$ generated by
$\Id$ and the operators $L_{a}\colon x\mapsto a\bullet x$ is, for semisimple
$J$, a semisimple \emph{associative} algebra; hence every $J$-module is
completely reducible and the cohomology splits. This is the engine behind
\Cref{thm:whitehead}.

% PROV: FarautKoranyi1994 via agent-A/notes/ingestion-results-2026-06-01.json (item 5) -- Aut(J) compact for Euclidean J, closed in O(J, trace form). Extraction-level provenance; local primary text still pending.
\begin{proposition}[Compactness of the automorphism group, \status{(source-extracted; pending primary byte-check)}]\label{prop:aut-compact}
For a finite-dimensional Euclidean (formally real) Jordan algebra $J$ the
automorphism group $\Aut(J)$ is compact: it is a closed subgroup of the
orthogonal group of the trace inner product, and every automorphism is an
isometry of that inner product. Consequently normalized Haar averaging over
$\Aut(J)$ is available, producing invariant inner products and invariant
tensors.
\end{proposition}

\noindent \Cref{prop:aut-compact} is the analytic input recorded in
\cite{FarautKoranyi1994}: Jordan automorphisms of a Euclidean Jordan algebra
preserve the trace form, so $\Aut(J)$ is a closed, hence compact, subgroup of an
orthogonal group. This is the Jordan counterpart of the compact unitary group
that Kitaev averages over.

\smallskip\noindent\textbf{Agent B review provenance flag.}
\Cref{thm:whitehead,prop:aut-compact} are registered in the provenance ledger
through the local ingestion file
\path{agent-A/notes/ingestion-results-2026-06-01.json}.
The Chu--Russo primary source is present locally as a PDF, but the Whitehead
quotes have not yet been independently byte-matched against extracted primary
text; Faraut--Koranyi compactness is likewise extraction-level provenance in
this draft.

\subsection{The dimension-free obstruction and candidate routes}

We now isolate the single quantitative difficulty in step (d), state the
obstruction precisely, and describe the resolution proposed in the project's
internal notes. This subsection is \textbf{original} to the project and is
explicitly a \textbf{programme}, marked \status{(open)}; it is being tested
numerically.

\paragraph{The cochain setup.}
Write $g(b,b'):=v(b\bullet b')-v(b)\bullet v(b')$ for the multiplicativity
defect of the inclusion $v\colon\mathcal{B}\to\mathcal{A}$ from step (c); it is
symmetric with $\norm{g}\le\delta$. Linearizing the Jordan identity (which
$\mathcal{B}$ satisfies exactly and $\mathcal{A}$ only up to $\eps$) shows that
$g$ is an \emph{$\eps$-approximate Jordan $2$-cocycle}: the cocycle identity holds
up to $O(\eps)\norm{\cdot}^{3}$. By \Cref{thm:whitehead} (with $J=\mathcal{B}$)
an exact cocycle would be an exact coboundary, $g=d^{1}h$ for a linear
$h\colon\mathcal{B}\to\mathcal{A}$, where
$(d^{1}h)(b,b')=v(b)\bullet h(b')+h(b)\bullet v(b')-h(b\bullet b')$. Setting
$\tilde v:=v-h$ replaces the defect by
$\tilde g=g-d^{1}h+O(\delta^{2})=O(\delta^{2}+K\eps)$, where $K=\norm{s}$ is the
norm of the splitting operator $s$ that produces $h=sg$. Iterating the resulting
Newton step $\delta\mapsto O(\delta^{2}+K\eps)$ drives any sufficiently small
$\delta$ down to $O(K\eps)$. The theorem therefore reduces to one quantitative
question: \emph{is $K$ dimension-free?}

\paragraph{The obstruction.}
The splitting $s$ is built from a \emph{diagonal} (Casimir) tensor
$D=\sum_{j}u_{j}\otimes u_{j}'$ of the multiplication algebra $M(\mathcal{B})$;
the natural candidate is $D=\sum_{i}e_{i}\otimes e^{i}$ formed from a pair of
trace-form-dual bases. The relevant norm is the projective tensor norm
$\norm{\cdot}_{\wedge}$, an infimum over representations. Evaluated on the
\emph{basis representation}, this naive Casimir has
\[
   \norm{D}_{\wedge}\;\sim\;\dim\mathcal{B},
\]
so the resulting $K$ grows with the dimension and the Newton iteration only
converges when $\eps<c/\dim\mathcal{B}$. This is exactly the dimension-dependent
bound Kitaev encounters for the naive associative diagonal: he notes that
``naive constructions of the Haar measure (or just the diagonal) in the
$\eps$-associative setting have error bounds proportional to
$n=\dim\mathcal{A}$'' \cite{Kitaev2024}.

\paragraph{Candidate averaging mechanism (original).}
The project note registered in \texttt{PROVENANCE.md} as \texttt{A-ER} proposes
that $\norm{\cdot}_{\wedge}$ should be estimated using a compact-group
representation rather than a basis representation. In the associative case
Kitaev's diagonal is the group average
\[
   D=\int_{U(n)}U^{\dagger}\otimes U\,dU=\tfrac1n\,\mathrm{SWAP},
\]
\cite{Kitaev2024}. The basis representation of $\tfrac1n\,\mathrm{SWAP}$ gives
$\norm{D}_{\wedge}\le\tfrac1n\cdot n^{2}=n$, but the \emph{averaging}
representation gives
\[
   \norm{D}_{\wedge}\le\int\norm{U^{\dagger}}_{\mathrm{op}}\,
   \norm{U}_{\mathrm{op}}\,dU=\int 1\,dU=1 .
\]
It is the \emph{same} element; only the compact-group average exhibits the
$O(1)$ projective norm. The proposed Jordan analogue tries to imitate this by
using, not $\mathcal{B}\otimes\mathcal{B}$ directly, but the
\emph{separability idempotent}
$e=\sum_{i}u_{i}\otimes u_{i}'\in M(\mathcal{B})\otimes M(\mathcal{B})^{\mathrm{op}}$
of the semisimple associative algebra $M(\mathcal{B})$. In the trace/Frobenius
model one has a compact unitary group $U(M(\mathcal{B}))$ and a
dimension-free-looking average of the form
\[
   e=\int_{U(M(\mathcal{B}))}u^{\dagger}\otimes u\,du,
   \qquad
   \norm{e}_{\wedge}\le\int\norm{u^{\dagger}}\,\norm{u}\,du=1,
\]
which explains the favourable Frobenius-norm numerics below. This is not yet a
proof in the order-unit/operator norm required by \Cref{thm:jordan-structure}:
the unitary group of $M(\mathcal{B})$ preserves the trace norm, while the
order-unit norm is controlled by the generally smaller Jordan automorphism
group. The open task is to turn this candidate average, or an incremental
replacement for it, into a contracting homotopy $s$ with order-unit-norm bound
$\norm{s}=K=O(1)$.

% PROV: ORIGINAL agent-A/theory/01-error-reduction.md (§3) -- candidate group-averaged separability idempotent of M(B); order-unit dimension-free K remains open.
\begin{openproblem}[Dimension-free error reduction]\label{prop:error-reduction}
\status{(open; programme target).} If the compact-group-averaged
separability idempotent of the multiplication algebra $M(\mathcal{B})$ can be
converted into a contracting homotopy $s$ for Jordan $2$-cocycles with
order-unit-norm bound $\norm{s}=K=O(1)$, independent of $\dim\mathcal{B}$, then
the Newton iteration $\delta\mapsto O(\delta^{2}+K\eps)$ of step (d) terminates
at $\delta=O(\eps)$ with a dimension-free constant, which is what
\Cref{thm:jordan-structure} requires.
\end{openproblem}

The status of \Cref{prop:error-reduction} is the following. The two algebraic
inputs --- vanishing cohomology (\Cref{thm:whitehead}) and a compact averaging
group (\Cref{prop:aut-compact}) --- are cited facts. The mechanism that makes
the constant dimension-free is established for the associative diagonal by
Kitaev; its transfer to the Jordan setting has been examined numerically, with a
result that is informative and must be reported precisely.

\paragraph{Numerical evidence \status{(computational evidence; not theorem)}, and the precise remaining gap.}
We computed the smallest nonzero singular value of $d^{1}$, hence
$\norm{s}=1/\sigma_{\min}(d^{1})$, for $J=H_{n}(\C)$, $J=H_{n}(\R)$, and the
spin factors $V_{n}$, in two norms on the cochain spaces. (The computation is
validated by $\dim\Ker d^{1}=\dim\mathrm{Der}(J)$: equal to $n^{2}-1$ for
$H_{n}(\C)$ and $n(n-1)/2$ for $V_{n}$, both matched exactly.) The outcome
splits along the choice of norm:
\begin{itemize}[leftmargin=*]
\item \emph{Trace (Frobenius) inner-product norm:} $\norm{s}$ is
\textbf{bounded --- in fact $\le 1$ and decreasing} in every family
($H_{n}(\C)$: $0.91,0.78,0.69,0.62,0.57$ for $n=2,\dots,6$; $V_{n}$:
convergent to $\approx0.62$). In this norm the dimension-free splitting exists.
\item \emph{Order-unit (operator) norm --- the norm \Cref{thm:jordan-structure}
actually requires:} certified lower bounds on $\norm{s}_{\mathrm{op}}$ for
$H_{n}(\C)$ grow only \emph{slowly} ($1.4\to1.8$ for $n=2,\dots,6$, with
$\norm{s}_{\mathrm{op}}/\sqrt n$ and $\norm{s}_{\mathrm{op}}/n$ both decreasing),
while a crude upper bound grows like $n$. \textbf{Boundedness in the operator
norm is therefore unresolved} at this range.
\end{itemize}
The discrepancy is structural and locates the gap exactly. The Frobenius bound
comes from averaging over the unitary group of the multiplication algebra
$M(J)$, whose elements are Frobenius isometries. The order-unit norm, however,
is preserved only by the Jordan automorphism group $\Aut(J)$, which is
\emph{strictly smaller} than $U(M(J))$. The $O(1)$ projective-norm
representation of the separability idempotent uses the larger group, so it does
not by itself give an operator-norm bound. Closing the gap requires one of:
(i) an $\Aut(J)$-only construction of the splitting carrying an operator-norm
bound; (ii) Kitaev's incremental route, applying the splitting only to a
controlled algebra $\mathcal{B}$ built one step at a time, so that the dimension
never enters; or (iii) a formulation in the trace norm with an explicit,
possibly dimension-dependent, conversion. Independently of the norm, one must
still write the explicit contraction $h=sg$, verify $d^{1}h=g+O(\eps)$ on
$\eps$-approximate cocycles, and absorb the approximate-module errors (the
module $\mathcal{A}$ is itself only an $\eps$-JB algebra). Until these are
complete, \Cref{thm:jordan-structure} and \Cref{prop:error-reduction} remain
\status{(open)}. The numerical detail is recorded in
\texttt{agent-A/experiments/jordan-coboundary/REPORT.md}.
